This chapter reviews the research context in which L2M is situated. The
review is organized in four parts: general deep-learning approaches to music
generation, supervised lyrics-to-melody systems, text-conditioned audio
generation, and the recent application of LLMs to musical tasks. The chapter
closes with a gap analysis that motivates the design of L2M.

\section{Deep Learning for Music Generation}
\label{sec:lr-dl-music}

Deep learning approaches to symbolic music generation have employed
recurrent networks, GANs, VAEs, and Transformer architectures to model the
hierarchical temporal structure of music; Briot et
al.~\cite{briot20deeplearning} provide a comprehensive survey of these
foundational methods. Sequence-to-sequence learning
\cite{sutskever14seq2seq} first made it practical to map one symbolic
sequence to another of different length, forming the algorithmic basis for
lyric-conditioned melody generation. The Music Transformer
\cite{huang18musictransformer} subsequently demonstrated that self-attention
with relative position representations can capture long-range musical
structure such as repetition and motif development, establishing the
Transformer \cite{vaswani17attention} as the dominant architecture for
symbolic music modeling.

\section{Lyrics-to-Melody Generation}
\label{sec:lr-l2m}

\subsection{Early Sequence-to-Sequence Approaches}

Bao et al.~\cite{bao19neural} introduced one of the earliest end-to-end
neural frameworks for melody composition from lyrics, using syllable-level
alignment on pop music corpora within a Seq2Seq architecture. The work
demonstrated feasibility but required a large paired lyrics--melody corpus,
and generated melodies frequently violated basic music-theoretic
expectations. SongMASS \cite{sheng21songmass} improved data efficiency by
pre-training separate lyric and melody encoders with masked sequence-to-sequence
objectives and enforcing alignment constraints during decoding.

\subsection{Music-Theory-Constrained and Template-Based Systems}

ReLyMe \cite{zhang22relyme} addressed a core weakness of purely data-driven
models --- weak semantic and musical coherence --- by explicitly
incorporating lyric--melody relationship principles from music theory (tone,
rhythm, and structure relationships) as decoding-time constraints.
TeleMelody \cite{ju22telemelody} decomposed generation into a two-stage,
template-based process (lyric-to-template and template-to-melody), improving
controllability and reducing the dependence on paired data. ROC
\cite{lv22roc} went further and reformulated generation as retrieval and
re-composition over a database of existing melodic fragments. These systems
collectively demonstrate the value of decomposition and of explicit musical
knowledge --- two ideas that L2M adopts in a radically simplified, prompt-based
form.

\subsection{Controllable and Alignment-Focused Systems}

Zhang et al.~\cite{zhang23controllable} introduced style modulation via
reference embeddings and memory fusion, enabling user-directed genre and
mood control and marking the shift from pure generation toward interactive
music creation. Reddy et al.~\cite{reddy23attention} addressed incomplete or
partial lyric inputs by applying cross-attention between text and melody
sequences, improving robustness in real-world scenarios. More recent
conditional-Transformer systems provide fine-grained melodic control:
CSL-L2M \cite{chai24csll2m} conditions on song-level lyric and musical
attributes, while SongGLM \cite{yu24songglm} introduces 2D alignment
encoding with multi-task pre-training for harmonized lyric--melody
relationships. On the representation side, Wang et
al.~\cite{wang25melodylyrics} learn a shared melody--lyrics embedding space
with a contrastive alignment loss, demonstrating the importance of strong
cross-modal representations for both generation and retrieval.

\subsection{LLM-Based Song Composition}

SongComposer \cite{ding24songcomposer} is the closest recent work to this
thesis: a music-specialized LLM that generates lyrics and melodies jointly
using a tuple representation for word-level alignment. It achieves high
quality but requires extensive domain-specific fine-tuning on large paired
song datasets --- precisely the requirement L2M eliminates through zero-shot
prompting of a general-purpose model. ChatMusician \cite{yuan24chatmusician}
similarly extends an open LLM with intrinsic musical abilities via continued
pre-training on ABC notation, again at substantial training cost.

\section{Text-to-Audio Music Generation}
\label{sec:lr-tta}

A parallel research line generates \emph{audio} rather than symbolic
notation. MusicLM \cite{agostinelli23musiclm} casts text-conditioned music
generation as hierarchical sequence-to-sequence modeling over audio tokens,
and MusicGen \cite{copet23musicgen} showed that a single-stage language
model over compressed audio representations suffices for controllable,
high-quality generation. These systems produce rich, fully arranged music,
but they differ fundamentally from the task addressed here: they do not
enforce syllabic alignment with lyrical text, and their audio output cannot
be edited as notation by a musician. Recent comprehensive reviews of
AI-enabled text-to-music generation \cite{li25review} classify systems along
these symbolic-versus-audio axes and identify persistent challenges:
long-term coherence, data scarcity, and limited controllability ---
challenges that motivate symbolic, interpretable approaches such as L2M.

Related transformer-based work has also targeted the inverse and adjacent
tasks, such as lyric text generation with T5-style models
\cite{mediakov24lyrics} and LSTM-based lyric modeling \cite{gill20lyric},
which established sequential baselines for the language side of song
composition.

\section{Gap Analysis}
\label{sec:lr-gap}

\ref{tab:gap-analysis} summarizes the limitations shared by existing
lyrics-to-melody systems.

\begin{table}[ht]
\centering
\caption[Gap analysis of existing lyrics-to-melody systems]{Gap analysis of
existing lyrics-to-melody systems and the corresponding L2M design response.}
\label{tab:gap-analysis}
\begin{tabular}{p{6.3cm} p{5.8cm}}
\toprule
\textbf{Limitation in prior work} & \textbf{L2M design response} \\
\midrule
Require large paired lyrics--melody training corpora & Zero-shot LLM
prompting; no training data \\
No graceful degradation when the model fails at inference & Deterministic,
music-theory-grounded fallback at each LLM stage \\
Single output format (typically MIDI only) & MIDI, MusicXML, WAV, and MP3
export \\
No explicit emotion-to-key grounding & Explicit emotion-to-musical-attribute
mapping \\
Not available as deployable open-source software & Documented Python package
with CLI and API \\
\bottomrule
\end{tabular}
\end{table}

No prior system, to the best of our knowledge, combines zero-shot LLM
generation with a deterministic reliability layer and multi-format notation
export in a single deployable package. This combination defines the research
space that the remainder of this thesis develops.
